\chapter{Operadores Derivados: NAND, NOR, XOR y XNOR}
Las ecuaciones obtenidas a partir de las Leyes de De Morgan demuestran que las operaciones básicas $+$ y $\cdot$ pueden ser expresadas íntegramente en términos de la negación de su operación dual. Esto motiva la definición de varios operadores lógicos fundamentales en sistemas digitales. Dos de ellos (NAND y NOR) son de gran relevancia por ser funcionalmente completos por sí solos, mientras que otros dos (XOR y XNOR) son esenciales para funciones aritméticas y de comprobación de paridad:

\begin{definicion}{Operador NAND (Barra de Sheffer)}{nand}
Denotado clásicamente con una flecha hacia arriba ($\uparrow$), se define como la negación del ínfimo.
\[ a \uparrow b \triangleq \overline{(a \cdot b)} = \overline{a} + \overline{b} \]
\end{definicion}

\begin{definicion}{Operador NOR (Flecha de Peirce)}{nor}
Denotado con una flecha hacia abajo ($\downarrow$), se define como la negación del supremo.
\[ a \downarrow b \triangleq \overline{(a + b)} = \overline{a} \cdot \overline{b} \]
\end{definicion}

\begin{definicion}{Operador XOR (O-exclusiva)}{xor}
Denotado con el símbolo de suma exclusiva ($\oplus$), evalúa a $1$ cuando exactamente uno de los operandos es $1$ y el otro $0$.
\[ a \oplus b \triangleq (a \cdot \overline{b}) + (\overline{a} \cdot b) \]
\end{definicion}

\begin{definicion}{Operador XNOR (No-O-exclusiva o Equivalencia)}{xnor}
Denotado frecuentemente con $\odot$ o $\leftrightarrow$, es la negación de la operación XOR y evalúa a $1$ cuando ambos operandos son idénticos.
\[ a \odot b \triangleq \overline{(a \oplus b)} = (a \cdot b) + (\overline{a} \cdot \overline{b}) \]
\end{definicion}

\begin{definicion}{Generalización a $n$ variables de NAND y NOR}{gen_nand_nor}
A diferencia de los operadores $+$, $\cdot$ y $\oplus$, los operadores NAND ($\uparrow$) y NOR ($\downarrow$) \textbf{no son asociativos}. Sin embargo, debido a su inmensa importancia práctica en la construcción de circuitos digitales, se define convencionalmente su generalización a $n$ variables como la negación de la conjunción o disyunción múltiple, respectivamente:
\[ \text{NAND}(x_1, x_2, \dots, x_n) \triangleq \overline{\left( \prod_{i=1}^n x_i \right)} \]
\[ \text{NOR}(x_1, x_2, \dots, x_n) \triangleq \overline{\left( \sum_{i=1}^n x_i \right)} \]
\end{definicion}

\section{Comportamiento de los Operadores Derivados}
Los operadores NAND ($\uparrow$) y NOR ($\downarrow$) presentan una serie de propiedades algebraicas particulares. Al ser operadores funcionalmente completos, permiten expresar cualquier otra operación booleana utilizando exclusivamente uno de ellos.

\begin{teorema}{Idempotencia cruzada (Generación de NOT)}{idemp_cruzada}
Operar un elemento consigo mismo usando NAND o NOR equivale a su complemento (negación):
\begin{align*}
    \forall a \in \mathbb{B}, \quad a \uparrow a &= \overline{a} \\
    \forall a \in \mathbb{B}, \quad a \downarrow a &= \overline{a}
\end{align*}
\end{teorema}
\desplegable{proof_idemp_cruzada}{
\begin{proof}[Demostración]
Para la operación NAND:
\begin{align*}
    a \uparrow a &\triangleq \overline{(a \cdot a)} & (\text{Definicion de NAND}) \\
                 &= \overline{a} & (Idemp_\cdot)
\end{align*}
Para la operación NOR:
\begin{align*}
    a \downarrow a &\triangleq \overline{(a + a)} & (\text{Definicion de NOR}) \\
                   &= \overline{a} & (Idemp_+)
\end{align*}
\end{proof}
}

\begin{teorema}{Generación del Ínfimo y Supremo (AND y OR)}{gen_inf_sup}
A partir de la propiedad anterior, podemos recuperar las operaciones básicas anidando las puertas consigo mismas:
\begin{align*}
    \forall a, b \in \mathbb{B}, \quad a \cdot b &= \overline{(a \uparrow b)} = (a \uparrow b) \uparrow (a \uparrow b) \\
    \forall a, b \in \mathbb{B}, \quad a + b &= \overline{(a \downarrow b)} = (a \downarrow b) \downarrow (a \downarrow b)
\end{align*}
\end{teorema}
\desplegable{proof_gen_inf_sup}{
\begin{proof}[Demostración]
Para la generación del ínfimo (AND):
\begin{align*}
    \overline{(a \uparrow b)} &\triangleq \overline{(\neg(a \cdot b))} & (\text{Definicion de NAND}) \\
                       &= a \cdot b & (Involucion)
\end{align*}
Además, por la idempotencia cruzada demostrada anteriormente, $x \uparrow x = \overline{x}$, por tanto:
\[ \overline{(a \uparrow b)} = (a \uparrow b) \uparrow (a \uparrow b) \]
La demostración para la generación del supremo (OR) es idéntica por dualidad.
\end{proof}
}

\begin{teorema}{Generación cruzada (Leyes de De Morgan para NAND/NOR)}{gen_cruzada}
Podemos generar la operación opuesta (supremo desde NAND, e ínfimo desde NOR) negando previamente las entradas:
\begin{align*}
    \forall a, b \in \mathbb{B}, \quad a + b &= (\overline{a}) \uparrow (\overline{b}) = (a \uparrow a) \uparrow (b \uparrow b) \\
    \forall a, b \in \mathbb{B}, \quad a \cdot b &= (\overline{a}) \downarrow (\overline{b}) = (a \downarrow a) \downarrow (b \downarrow b)
\end{align*}
\end{teorema}
\desplegable{proof_gen_cruzada}{
\begin{proof}[Demostración]
Para el supremo (OR):
\begin{align*}
    (\overline{a}) \uparrow (\overline{b}) &\triangleq \overline{(\neg a \cdot \neg b)} & (\text{Definicion de NAND}) \\
                               &= \overline{(\neg (a + b))} & (Mor_+) \\
                               &= a + b & (Involucion)
\end{align*}
La demostración para el ínfimo (AND) sigue los mismos pasos de manera dual, aplicando $Mor_\cdot$.
\end{proof}
}

\begin{teorema}{Conmutatividad}{conmut_deriv}
Al igual que sus operaciones base, ambos operadores son perfectamente conmutativos:
\begin{align*}
    \forall a, b \in \mathbb{B}, \quad a \uparrow b &= b \uparrow a \\
    \forall a, b \in \mathbb{B}, \quad a \downarrow b &= b \downarrow a
\end{align*}
\end{teorema}
\desplegable{proof_conmut_deriv}{
\begin{proof}[Demostración]
Para la operación NAND:
\begin{align*}
    a \uparrow b &\triangleq \overline{(a \cdot b)} & (\text{Definicion de NAND}) \\
                 &= \overline{(b \cdot a)} & (Comm_\cdot) \\
                 &\triangleq b \uparrow a & (\text{Definicion de NAND})
\end{align*}
Para la operación NOR, es análogo aplicando $Comm_+$.
\end{proof}
}

\begin{teorema}{Inexistencia de Elemento Neutro}{no_neutro}
No existe ningún elemento neutro para las operaciones NAND ni NOR en un álgebra de Boole general.
\end{teorema}
\desplegable{proof_no_neutro}{
\begin{proof}[Demostración]
Si existiera un neutro $e$ para la operación NAND, debería cumplirse que $\forall a, a \uparrow e = a$, es decir, $\overline{(a \cdot e)} = a$. Si probamos con $e=1$, obtenemos $\overline{a} = a$, lo cual obliga al colapso en un álgebra trivial. Si probamos con $e=0$, obtenemos $\overline{0} = a \implies 1 = a$, lo cual obviamente no se cumple para cualquier elemento $a$. Lo mismo aplica a la operación NOR.
\end{proof}
}

\begin{teorema}{Comportamiento con las constantes (Fijación y Absorción)}{constantes_deriv}
Fijar una constante específica en uno de los operandos genera directamente la negación, mientras que usar la constante opuesta actúa como un pseudo-elemento absorbente (devolviendo un valor constante inalterable por $a$):
\begin{align*}
    \text{Inversión: } & a \uparrow 1 = \overline{a} \qquad & a \downarrow 0 &= \overline{a} \\
    \text{Absorción: } & a \uparrow 0 = 1 \qquad & a \downarrow 1 &= 0
\end{align*}
\end{teorema}
\desplegable{proof_constantes_deriv}{
\begin{proof}[Demostración]
Para la inversión:
\begin{align*}
    a \uparrow 1 &= \overline{(a \cdot 1)} = \overline{a} & (ElemNeu_\cdot) \\
    a \downarrow 0 &= \overline{(a + 0)} = \overline{a} & (ElemNeu_+)
\end{align*}
Para la absorción:
\begin{align*}
    a \uparrow 0 &= \overline{(a \cdot 0)} = \overline{0} = 1 & (Abs_0) \\
    a \downarrow 1 &= \overline{(a + 1)} = \overline{1} = 0 & (Abs_1)
\end{align*}
\end{proof}
}

\begin{teorema}{Ausencia de Asociatividad}{no_asoc_deriv}
A diferencia del supremo ($+$) y el ínfimo ($\cdot$), las operaciones NAND y NOR son positivamente NO asociativas:
\begin{align*}
    (a \uparrow b) \uparrow c &\neq a \uparrow (b \uparrow c) \\
    (a \downarrow b) \downarrow c &\neq a \downarrow (b \downarrow c)
\end{align*}
\end{teorema}
\desplegable{proof_no_asoc_deriv}{
\begin{proof}[Demostración]
Desarrollando el lado izquierdo para NAND:
\[ (a \uparrow b) \uparrow c = \overline{( (\neg (a \cdot b)) \cdot c )} = (a \cdot b) + \overline{c} \]
Desarrollando el lado derecho:
\[ a \uparrow (b \uparrow c) = \overline{( a \cdot (\neg (b \cdot c)) )} = \overline{a} + (b \cdot c) \]
Resulta evidente que $(a \cdot b) + \overline{c} \neq \overline{a} + (b \cdot c)$ para combinaciones arbitrarias de variables. El mismo razonamiento aplica de manera estricta y análoga para demostrar la carencia de asociatividad en la operación NOR ($\downarrow$).
\end{proof}
}

\begin{definicion}{NAND y NOR de $n$ entradas}{def_n_entradas}
Dado que las puertas NAND y NOR físicas a menudo tienen más de dos entradas, se definen algebraicamente para múltiples entradas como la negación de la conjunción o disyunción de todas ellas:
\begin{align*}
    \uparrow(x_1, x_2, \dots, x_n) &\triangleq \overline{\left( \prod_{i=1}^n x_i \right)} \\
    \downarrow(x_1, x_2, \dots, x_n) &\triangleq \overline{\left( \sum_{i=1}^n x_i \right)}
\end{align*}
En particular, para el caso de 3 entradas que estudiaremos a continuación:
\begin{align*}
    \uparrow(a,b,c) &\triangleq \overline{(a \cdot b \cdot c)} \\
    \downarrow(a,b,c) &\triangleq \overline{(a + b + c)}
\end{align*}
\end{definicion}

\begin{teorema}{NAND/NOR múltiple vs agrupación binaria}{multiple_vs_binaria}
Como consecuencia directa de su falta de asociatividad, una operación NAND o NOR de 3 entradas no es equivalente a la agrupación secuencial en cascada de operaciones de 2 entradas:
\begin{align*}
    \uparrow(a,b,c) &\neq (a \uparrow b) \uparrow c \qquad \text{y} \qquad \uparrow(a,b,c) \neq a \uparrow (b \uparrow c) \\
    \downarrow(a,b,c) &\neq (a \downarrow b) \downarrow c \qquad \text{y} \qquad \downarrow(a,b,c) \neq a \downarrow (b \downarrow c)
\end{align*}
\end{teorema}
\desplegable{proof_multiple_vs_binaria}{
\begin{proof}[Demostración]
Para la NAND de 3 entradas tenemos, por definición y leyes de De Morgan:
\[ \uparrow(a,b,c) \triangleq \overline{(a \cdot b \cdot c)} = \overline{a} + \overline{b} + \overline{c} \]
Sin embargo, la agrupación de dos en dos evaluada anteriormente daba:
\[ (a \uparrow b) \uparrow c = (a \cdot b) + \overline{c} \]
Evidentemente, $\overline{a} + \overline{b} + \overline{c} \neq (a \cdot b) + \overline{c}$. Lo mismo aplica a las agrupaciones derechas y a las operaciones NOR equivalentes.
\end{proof}
}

\begin{teorema}{Extensión del Principio de Dualidad (NAND y NOR)}{dualidad_nand_nor}
La inclusión de los operadores derivados expande el Principio de Dualidad establecido en los postulados iniciales. La expresión dual de cualquier teorema o identidad que contenga operaciones NAND o NOR se obtiene intercambiando los operadores $\uparrow$ y $\downarrow$ (además de los ya conocidos $+ \leftrightarrow \cdot$ y $0 \leftrightarrow 1$).
\end{teorema}
\subsection{Comportamiento de los Operadores XOR y XNOR}
A diferencia de los operadores NAND y NOR, que destacan por su universalidad funcional pero carecen de propiedades algebraicas deseables (como asociatividad o elemento neutro), los operadores XOR ($\oplus$) y XNOR ($\odot$) exhiben una rica estructura algebraica. A continuación demostraremos estas propiedades.

\begin{teorema}{Conmutatividad}{conmut_xor}
Ambos operadores son conmutativos:
\[ a \oplus b = b \oplus a \qquad \text{y} \qquad a \odot b = b \odot a \]
\end{teorema}
\desplegable{proof_conmut_xor}{
\begin{proof}[Demostración]
Para XOR:
\begin{align*}
    a \oplus b &\triangleq (a \cdot \overline{b}) + (\overline{a} \cdot b) & (\text{Definición de XOR}) \\
               &= (\overline{a} \cdot b) + (a \cdot \overline{b}) & (Comm_+) \\
               &= (b \cdot \overline{a}) + (\overline{b} \cdot a) & (Comm_\cdot) \\
               &\triangleq b \oplus a & (\text{Definición de XOR})
\end{align*}
La demostración para XNOR sigue pasos idénticos aprovechando la conmutatividad del ínfimo y el supremo.
\end{proof}
}

\begin{teorema}{Elementos Neutros e Inversores}{neutro_xor}
El elemento $0$ actúa como neutro para la XOR, y $1$ actúa como inversor. De manera dual, $1$ es el neutro de la XNOR, y $0$ actúa como inversor:
\begin{align*}
    a \oplus 0 &= a & a \oplus 1 &= \overline{a} \\
    a \odot 1 &= a & a \odot 0 &= \overline{a}
\end{align*}
\end{teorema}
\desplegable{proof_neutro_xor}{
\begin{proof}[Demostración para XOR]
Para el elemento $0$ (neutro):
\begin{align*}
    a \oplus 0 &\triangleq (a \cdot \overline{0}) + (\overline{a} \cdot 0) & (\text{Definición}) \\
                  &= (a \cdot 1) + 0 & (Comp_{inv} \text{ y } Fij_\cdot) \\
                  &= a + 0 & (ElemNeu_\cdot) \\
                  &= a & (ElemNeu_+)
\end{align*}
Para el elemento $1$ (inversor):
\begin{align*}
    a \oplus 1 &\triangleq (a \cdot \overline{1}) + (\overline{a} \cdot 1) & (\text{Definición}) \\
                  &= (a \cdot 0) + \overline{a} & (Comp_{inv} \text{ y } ElemNeu_\cdot) \\
                  &= 0 + \overline{a} & (Fij_\cdot) \\
                  &= \overline{a} & (ElemNeu_+)
\end{align*}
Las pruebas para XNOR son totalmente duales.
\end{proof}
}

\begin{teorema}{Elemento Inverso de sí mismo (Grupo Abeliano)}{idemp_nula_xor}
La combinación de un elemento consigo mismo produce una anulación (devuelve el neutro de la operación correspondiente), actuando cada elemento como su propio inverso:
\[ a \oplus a = 0 \qquad \text{y} \qquad a \odot a = 1 \]
\end{teorema}
\desplegable{proof_idemp_nula_xor}{
\begin{proof}[Demostración]
Para XOR:
\begin{align*}
    a \oplus a &\triangleq (a \cdot \overline{a}) + (\overline{a} \cdot a) \\
               &= 0 + 0 & (Comp_\cdot \text{ y } Comm_\cdot) \\
               &= 0 & (Idemp_+)
\end{align*}
Para XNOR:
\begin{align*}
    a \odot a &\triangleq (a \cdot a) + (\overline{a} \cdot \overline{a}) \\
              &= a + \overline{a} & (Idemp_\cdot) \\
              &= 1 & (Comp_+)
\end{align*}
\end{proof}
}

\begin{teorema}{Propiedades de Negación}{negacion_xor}
Negar cualquiera de las entradas de forma independiente equivale a negar la operación completa, lo que a su vez alterna entre XOR y XNOR:
\begin{align*}
    \overline{(a \oplus b)} &= \overline{a} \oplus b = a \oplus \overline{b} = a \odot b \\
    \overline{(a \odot b)} &= \overline{a} \odot b = a \odot \overline{b} = a \oplus b
\end{align*}
\end{teorema}
\desplegable{proof_negacion_xor}{
\begin{proof}[Demostración de $a \oplus \overline{b} = \overline{(a \oplus b)}$]
\begin{align*}
    a \oplus \overline{b} &\triangleq (a \cdot \overline{(\neg b)}) + (\overline{a} \cdot \overline{b}) \\
                    &= (a \cdot b) + (\overline{a} \cdot \overline{b}) & (Involucion) \\
                    &\triangleq a \odot b & (\text{Definición de XNOR})
\end{align*}
Sabiendo por definición que $a \odot b \triangleq \overline{(a \oplus b)}$, se concluye de forma inmediata que $a \oplus \overline{b} = \overline{(a \oplus b)}$.
\end{proof}
}

\begin{teorema}{Asociatividad}{asoc_xor}
Tanto XOR como XNOR son operadores algebraicamente asociativos:
\[ (a \oplus b) \oplus c = a \oplus (b \oplus c) \qquad \text{y} \qquad (a \odot b) \odot c = a \odot (b \odot c) \]
\end{teorema}
\desplegable{proof_asoc_xor}{
\begin{proof}[Demostración para XOR]
Primero evaluaremos el miembro izquierdo $(a \oplus b) \oplus c$. Llamaremos $X = a \oplus b$.
\begin{align*}
    X \oplus c &\triangleq (X \cdot \overline{c}) + (\overline{X} \cdot c) \\
               &= \big( ((a \cdot \overline{b}) + (\overline{a} \cdot b)) \cdot \overline{c} \big) + \big( \overline{((a \cdot \neg b) + (\neg a \cdot b))} \cdot c \big) \\
               &= \big( ((a \cdot \overline{b}) + (\overline{a} \cdot b)) \cdot \overline{c} \big) + \big( (a \odot b) \cdot c \big) \quad (\text{Definición de XNOR}) \\
               &= \big( (a \cdot \overline{b} \cdot \overline{c}) + (\overline{a} \cdot b \cdot \overline{c}) \big) + \big( ((a \cdot b) + (\overline{a} \cdot \overline{b})) \cdot c \big) \quad (Dist_\cdot) \\
               &= (a \cdot \overline{b} \cdot \overline{c}) + (\overline{a} \cdot b \cdot \overline{c}) \\
               &\quad + (a \cdot b \cdot c) + (\overline{a} \cdot \overline{b} \cdot c) \quad (Dist_\cdot)
\end{align*}
Ahora evaluaremos el miembro derecho $a \oplus (b \oplus c)$. Llamaremos $Y = b \oplus c$.
\begin{align*}
    a \oplus Y &\triangleq (a \cdot \overline{Y}) + (\overline{a} \cdot Y) \\
               &= (a \cdot (b \odot c)) + (\overline{a} \cdot ((b \cdot \overline{c}) + (\overline{b} \cdot c))) \\
               &= (a \cdot ((b \cdot c) + (\overline{b} \cdot \overline{c}))) + (\overline{a} \cdot b \cdot \overline{c}) + (\overline{a} \cdot \overline{b} \cdot c) \\
               &= (a \cdot b \cdot c) + (a \cdot \overline{b} \cdot \overline{c}) \\
               &\quad + (\overline{a} \cdot b \cdot \overline{c}) + (\overline{a} \cdot \overline{b} \cdot c)
\end{align*}
Como podemos observar, ambas expansiones resultan exactamente en los mismos cuatro minitérminos. Reordenándolos por conmutatividad ($Comm_+$) demostramos que son idénticos.
\end{proof}
}

\begin{teorema}{Generalización n-aria (XOR y XNOR)}{gen_xor}
Dado que ambos operadores han demostrado ser asociativos y conmutativos, es posible omitir los paréntesis y generalizar la operación a un número arbitrario $n$ de variables. Se denotan mediante los operadores de sumatoria y productorio modificados:
\[ \bigoplus_{i=1}^{n} x_i = x_1 \oplus x_2 \oplus \dots \oplus x_n \]
\[ \bigodot_{i=1}^{n} x_i = x_1 \odot x_2 \odot \dots \odot x_n \]
\end{teorema}

\begin{teorema}{Distributividad (AND sobre XOR y OR sobre XNOR)}{dist_xor}
El producto (AND) se distribuye sobre la suma exclusiva (XOR), y dualmente, la suma (OR) se distribuye sobre la equivalencia (XNOR):
\[ a \cdot (b \oplus c) = (a \cdot b) \oplus (a \cdot c) \]
\[ a + (b \odot c) = (a + b) \odot (a + c) \]
\end{teorema}
\desplegable{proof_dist_xor}{
\begin{proof}[Demostración de AND sobre XOR]
Desarrollando el lado derecho (RHS):
\begin{align*}
    (a \cdot b) \oplus (a \cdot c) &\triangleq ((a \cdot b) \cdot \overline{(a \cdot c)}) + (\overline{(a \cdot b)} \cdot (a \cdot c)) \\
                                     &= (a \cdot b \cdot (\overline{a} + \overline{c})) + ((\overline{a} + \overline{b}) \cdot a \cdot c) \quad (Mor_\cdot) \\
                                     &= ((a \cdot b \cdot \overline{a}) + (a \cdot b \cdot \overline{c})) \\
                                     &\quad + ((a \cdot c \cdot \overline{a}) + (a \cdot c \cdot \overline{b})) \quad (Dist_\cdot) \\
                                     &= (0 + (a \cdot b \cdot \overline{c})) + (0 + (a \cdot \overline{b} \cdot c)) \quad (Comp_\cdot \text{ y } Fij_\cdot) \\
                                     &= (a \cdot b \cdot \overline{c}) + (a \cdot \overline{b} \cdot c) \quad (ElemNeu_+) \\
                                     &= a \cdot ((b \cdot \overline{c}) + (\overline{b} \cdot c)) \quad (Dist_\cdot \text{ a la inversa}) \\
                                     &\triangleq a \cdot (b \oplus c) \quad (\text{Def. XOR})
\end{align*}
Lo cual demuestra la igualdad. La demostración de la distributividad para XNOR es rigurosamente dual.
\end{proof}
}

\begin{teorema}{Extensión del Principio de Dualidad (XOR y XNOR)}{dualidad_xor_xnor}
De forma análoga a la relación entre NAND y NOR, los operadores XOR y XNOR son mutuamente duales. Para obtener la expresión dual de cualquier proposición que involucre estos operadores, se deben intercambiar $\oplus$ y $\odot$, manteniendo las reglas de dualidad estándar para los demás elementos y constantes.
\end{teorema}

